% FILE: 04_implementation_surface.tex
% PROJECT: experiments-core
% THESIS ROLE: benchmark-experiments-and-equation-verification

\section{Implementation Surface}
\label{sec:experiments-core-impl}

\subsection{Source Files}
\label{sec:experiments-core-impl-sources}

The implementation surface spans four categories of source files.

\textbf{Python validation scripts.} Three scripts constitute the executable equation
verification layer. \texttt{validate\_log\_fitness.py} encodes five Petri net token
replay scenarios---perfect conformance (fitness = 1.0), partial replay (fitness
= 0.3333), one-missing-token (fitness = 0.5), a combined multi-trace log (fitness
= 5/7), and a zero-consumed edge case---and asserts exact results using Python's
\texttt{assert} statement with no test framework dependency.
\texttt{validate\_ocpq\_refinement.py} encodes four OCPQ refinement relation cases:
reflexivity ($p \leq_L p$), extension ($p \leq_L c$ when $c$ extends $p$),
conflict ($p \not\leq_L c$ when values differ), and missing domain
($p \not\leq_L c$ when $\mathrm{dom}(p) \not\subseteq \mathrm{dom}(c)$).
\texttt{verify\_equations.py} is a single gate script that runs all nine cases
(five fitness, four OCPQ) and reports 100\% correctness or exits with a non-zero
status code.

\textbf{TypeScript visualizer layer.} \texttt{visualizer/alignment.js} implements
the A* alignment solver. \texttt{visualizer/blockchain.js} implements the SHA-256
cryptographic audit chain. \texttt{visualizer/drift-detector.js} implements the
EWMA drift detector. \texttt{visualizer/visualizer-validation.ts} validates ten WASM
boundary types at the TypeScript type level: \texttt{EvidenceTs}, \texttt{AdmissionTs},
\texttt{RefusalTs}, \texttt{ReceiptShapeTs}, \texttt{GraduationCandidateTs}, and
five supporting types. \texttt{visualizer/bindings.d.ts} declares the TypeScript
projections of the Rust typestate system.

\textbf{Benchmark reports.} \texttt{pm4py-comparison/matrix.md} records the
capability comparison matrix across DFG Discovery, Alpha Miner, Inductive Miner,
Heuristics Miner, Token Replay, Alignment Fitness, OCEL 2.0 handling, and ZKP
compatibility. \texttt{pm4py-comparison/zero-knowledge-benchmarks.md} records the
ZKP-specific comparison: wasm4pm completes 10,000-event ZKP conformance in $2^{18}$
cycles (approximately 142 seconds proof time, 220\,KB STARK) versus PM4Py timing
out at $2^{30}+$ cycles.

\textbf{Experiment sample documents.} Eighteen Markdown files with concrete JSON
schemas cover Petri net token replay, OCEL 2.0 lifecycle, cryptographic replay
receipts, OTel-BLAKE3 integration, raw laundering refusal, M\&A claim mappings,
paper-to-type-law fixture design, and decommission receipts.

\subsection{Commands}
\label{sec:experiments-core-impl-commands}

The primary commands are:

\begin{itemize}
  \item \texttt{python verify\_equations.py}---runs the unified equation gate; exit
    code 0 on all-pass, non-zero on any failure.
  \item \texttt{python validate\_log\_fitness.py}---runs the five Petri net token
    replay assertions.
  \item \texttt{python validate\_ocpq\_refinement.py}---runs the four OCPQ refinement
    assertions.
\end{itemize}

No build step is required for the Python scripts. The visualizer is a browser-based
application with no compilation step documented in the repository. The TypeScript
validation file (\texttt{visualizer-validation.ts}) requires a TypeScript compiler
to execute; no \texttt{tsconfig.json} or \texttt{package.json} is referenced in the
source evidence, so the precise compilation command is not reproducible from the
artifacts alone. [INTERPRETATION: The TypeScript validation is likely intended as a
type-check-only pass with \texttt{tsc --noEmit}.]

\subsection{Data and Ontology Surfaces}
\label{sec:experiments-core-impl-data}

The primary data surface is the OCEL 2.0 log schema, instantiated in concrete JSON
in the experiment sample documents. The ontology surface is the OCPQ query variable
domain---the refinement relation $p \leq_L c$ is defined over finite partial maps
from variable names to values. The OTel-BLAKE3 integration surface maps OpenTelemetry
span fields (TraceID, SpanID, ParentSpanID, SpanName, StartTime, EndTime,
InstructionCount) to BLAKE3 receipt chain inputs. No RDF or TTL ontology files
are present in the experiments directory; the ontology grounding for OCEL 2.0
is inherited from the broader process-intelligence research program. [FUTURE WORK:
A formal OWL ontology for the receipt-chain data model has not yet been manufactured.]

\subsection{Templates and Rendered Artifacts}
\label{sec:experiments-core-impl-templates}

The experiment sample Markdown files serve as templates for concrete JSON fixtures.
The \texttt{paper-fixture-design/spec.md} describes a specification for adversarial
fixture manufacturing using GAN-based trace generation with hyper-concurrent loops,
invisible prime branches, and 15\% missing events. However, no actual GAN
implementation or fixture files were materialized from this specification. The
\texttt{audit\_\_experiment\_completeness.md} is the rendered audit artifact that
confirms 18/18 experiments passed. The \texttt{checkpoint\_\_experiments\_complete.md}
is the rendered checkpoint that declares VERIFIED \& COMPLETE for the equation-verification
and experiment-coverage gates. The \texttt{EVIDENCE\_CHAIN\_TRACE.md} is the
rendered evidence-chain document that also identifies GAP\_001 as a blocking defect
in the full chain.
